\answer{ex:14:1}{(a)~分别拉长$0.17$\,s和$1.5$\,s。(b)~距离小于$11$\,cm时。}
\answer{ex:14:2}{只要$g\le\beta w_{q\mu}/w_{z\mu}=5$，无痛的触摸在任何$\mu$下都不通过；$g>6$时触摸$\mu=1$能通过：强烈的敏化使普通的触摸变成疼痛。}
\answer{ex:14:3}{(a)~$g^*=(1-k_sC)/(1-k_sA)$，当$k_sA<1$时稳定。(b)~无触摸时为$\nu\ge2$，有触摸时$\nu\ge1.7$即失控。}
\answer{ex:14:4}{$V(t)=-Pe^{-(T-t)/\tau_\gamma}$。(a)~$-Pe^{-T/\tau_\gamma}\approx-0.37P$。(b)~$+Pe^{-(T-t_e)/\tau_\gamma}\approx+0.61P$；它随$t_e$增大到$+P$，教会网络回避疼痛。}
\answer{ex:14:5}{$k_s=4$时不会自锁（$k_s\ge4.58$时才有第二个稳定状态）；$k_s=4.6$、$5$、$6$、$8$时$T_{\min}\approx4.35$、$1.40$、$0.65$、$0.32\,\tau_s$。}
\answer{ex:14:6}{$w_{q\nu}=4/9\approx0.444$，$q_0=2/9\approx0.222$，$w_{q\mu}=49/45\approx1.089$；无痛的触摸在$g=49/9\approx5.4$以内是安全的。}
\answer{ex:14:7}{(a)~当$\nu_c\ge q_0/w_{q\nu}$时阈值$\nu_c(g)=\theta/g$，否则为$(\theta+\beta q_0)/(g+\beta w_{q\nu})$；$g^*\approx2.45$、$1.0$、$0.25$。(b)~开放问题。}
